\section{Yang--Mills gradient flow and the small flow-time expansion}
\label{sec:3}
The Yang--Mills gradient flow defines a $D+1$~dimensional gauge
potential~$B(t,x)$ along a fictitious time~$t$, according to the flow equation
\begin{equation}
   \partial_tB_\mu(t,x)=D_\nu G_{\nu\mu}(t,x)
   +\alpha_0D_\mu\partial_\nu B_\nu(t,x),
\label{eq:(3.1)}
\end{equation}
where the $D+1$~dimensional field strength and the covariant derivative are
defined by
\begin{equation}
   G_{\mu\nu}(t,x)
   =\partial_\mu B_\nu(t,x)-\partial_\nu B_\mu(t,x)
   +[B_\mu(t,x),B_\nu(t,x)]
\label{eq:(3.2)}
\end{equation}
and
\begin{equation}
   D_\mu=\partial_\mu+[B_\mu,\cdot],
\label{eq:(3.3)}
\end{equation}
respectively. The initial condition for the flow is given by the
$D$~dimensional gauge potential in the previous section:
\begin{equation}
   B_\mu(t=0,x)=A_\mu(x).
\label{eq:(3.4)}
\end{equation}
In Eq.~\eqref{eq:(3.1)}, the last term is introduced to suppress the evolution
of the field along the direction of gauge degrees of freedom. Although this
term breaks the gauge symmetry, it does not affect the evolution of any
gauge-invariant operators~\cite{Luscher:2010iy}. Note that the mass dimension
of the flow time~$t$ is~$-2$.

Now, from the field strength extended to the $D+1$~dimension~\eqref{eq:(3.2)},
we define a $D+1$~dimensional analogue of the energy--momentum tensor by
\begin{equation}
   U_{\mu\nu}(t,x)\equiv
   G_{\mu\rho}^a(t,x)G_{\nu\rho}^a(t,x)
   -\frac{1}{4}\delta_{\mu\nu}G_{\rho\sigma}^a(t,x)G_{\rho\sigma}^a(t,x).
\label{eq:(3.5)}
\end{equation}
Although this is similar in form to the original energy--momentum
tensor~\eqref{eq:(2.4)}, it is not obvious a priori how this $D+1$~dimensional
object and Eq.~\eqref{eq:(2.4)} are related (or not). To find the relationship
between them is the principal task of the present paper. We also use the
density operator studied in~Ref.~\cite{Luscher:2010iy}:
\begin{equation}
   E(t,x)\equiv\frac{1}{4}G_{\mu\nu}^a(t,x)G_{\mu\nu}^a(t,x).
\label{eq:(3.6)}
\end{equation}

Now, as shown in~Ref.~\cite{Luscher:2011bx}, for~$t>0$, any correlation
function of~$B_\mu(t,x)$ is UV finite after standard renormalization in the
$D$~dimensional Yang--Mills theory. This property holds even for any local
products of~$B_\mu(t,x)$ such as~Eqs.~\eqref{eq:(3.5)} and~\eqref{eq:(3.6)}.
Also, for small flow times, a local product of~$B_\mu(t,x)$ can be regarded as
a local field in the $D$~dimensional sense because the flow
equation~\eqref{eq:(3.1)} is basically the diffusion equation along the
time~$t$ and the diffusion length in~$x$ is~$\sqrt{8t}$. These properties allow
us to express, as explained in~Sect.~8 of~Ref.~\cite{Luscher:2011bx},
$U_{\mu\nu}(t,x)$ and~$E(t,x)$ as an asymptotic series of $D$~dimensional
renormalized local operators with finite coefficients. Considering the gauge
invariance and the index structure, for $D=4$, we can write
\begin{equation}
   U_{\mu\nu}(t,x)
   =c_T(t)\left\{T_{\mu\nu}\right\}_R(x)
   +c_S(t)\delta_{\mu\nu}
   \left\{\frac{1}{4}F_{\rho\sigma}^aF_{\rho\sigma}^a\right\}_R(x)
   +O(t),
\label{eq:(3.7)}
\end{equation}
where abbreviated terms are the contributions of operators with a mass
dimension higher than or equal to~$6$. For Eq.~\eqref{eq:(3.6)}, we similarly
have
\begin{equation}
   E(t,x)
   =\left\langle E(t,x)\right\rangle
   +c_E(t)\left\{\frac{1}{4}F_{\rho\sigma}^aF_{\rho\sigma}^a\right\}_R(x)
   +O(t).
\label{eq:(3.8)}
\end{equation}
We note that, when the renormalized gauge coupling is fixed,
$U_{\mu\nu}(t,x)$~\eqref{eq:(3.5)} is traceless for~$D=4$,
\begin{equation}
   \delta_{\mu\nu}U_{\mu\nu}(t,x)
   =2\epsilon E(t,x)\xrightarrow{\epsilon\to0}0,
\label{eq:(3.9)}
\end{equation}
because $E(t,x)$~\eqref{eq:(3.6)} is finite~\cite{Luscher:2010iy} and does not
produce a $1/\epsilon$ singularity (this explains why there is no $c$~number
expectation value term in~Eq.~\eqref{eq:(3.7)}). Thus, considering the trace
part of~Eq.~\eqref{eq:(3.7)}, we see that the coefficients $c_T(t)$
and~$c_S(t)$ are not independent and are related by, for~$D=4$,
\begin{equation}
   c_S(t)=\frac{\beta}{2g^3}c_T(t),
\label{eq:(3.10)}
\end{equation}
because of the trace anomaly~\eqref{eq:(2.15)}.

By eliminating the renormalized action density from Eqs.~\eqref{eq:(3.7)}
and~\eqref{eq:(3.8)}, we have
\begin{equation}
   \left\{T_{\mu\nu}\right\}_R(x)
   =\frac{1}{c_T(t)}U_{\mu\nu}(t,x)
   -\frac{c_S(t)}{c_T(t)c_E(t)}\delta_{\mu\nu}
   \left[E(t,x)
   -\left\langle E(t,x)\right\rangle
   \right]
   +O(t).
\label{eq:(3.11)}
\end{equation}
This expression relates the energy--momentum tensor~\eqref{eq:(2.6)} and the
short flow-time behavior of gauge-invariant local products defined by the
gradient flow. Thus, once the coefficients are known, one can extract the
energy--momentum tensor from the $t\to0$ behavior of the combination on the
right-hand side.
